\documentclass[tikz,border=8pt]{standalone}
\usepackage{ctex}
\usepackage{amsmath}
\usetikzlibrary{arrows.meta,calc}
\begin{document}
\begin{tikzpicture}[>=Stealth]
  \tikzset{
    block/.style={rectangle, draw={rgb,255:red,74;green,144;blue,217}, fill=none,
      minimum width=1.82cm, minimum height=0.70cm, line width=1.2pt,
      font=\scriptsize\bfseries, align=center},
    wideblock/.style={rectangle, draw={rgb,255:red,74;green,144;blue,217}, fill=none,
      minimum width=2.15cm, minimum height=0.70cm, line width=1.2pt,
      font=\scriptsize\bfseries, align=center},
    sum/.style={circle, draw={rgb,255:red,92;green,184;blue,92}, fill=none,
      minimum size=0.44cm, line width=1.2pt,
      path picture={
        \draw[line width=0.8pt] (path picture bounding box.south west) --
          (path picture bounding box.north east);
        \draw[line width=0.8pt] (path picture bounding box.north west) --
          (path picture bounding box.south east);
      }},
    tap/.style={circle, draw={rgb,255:red,240;green,173;blue,78}, fill=none,
      minimum size=4.5pt, inner sep=0pt, line width=1.2pt},
    signal/.style={->, line width=1.2pt, draw={rgb,255:red,51;green,51;blue,51}},
    feedback/.style={->, line width=1.2pt, draw={rgb,255:red,51;green,51;blue,51}, dashed},
    plotbox/.style={rectangle, draw={rgb,255:red,198;green,208;blue,218}, fill=none,
      minimum width=1.70cm, minimum height=0.55cm, line width=0.8pt},
    labeltext/.style={font=\scriptsize, align=center}
  }

  \node[labeltext, font=\small\bfseries] at (6.7,2.65)
    {传感器噪声在一般闭环链路中的传播与滤波缓解};

  \newcommand{\drawrow}[4]{
    \node[labeltext, anchor=west, font=\scriptsize\bfseries] at (-0.35,#2+1.05) {#3};
    \node[block] (r#1) at (0,#2) {参考\\$r$};
    \node[sum] (s#1) at (2.05,#2) {};
    \node[block] (e#1) at (3.20,#2) {误差\\$e$};
    \node[block] (c#1) at (5.25,#2) {控制器};
    \node[wideblock] (a#1) at (7.65,#2) {执行机构};
    \node[wideblock] (p#1) at (10.15,#2) {对象输出};
    \node[tap] (tap#1) at (11.65,#2) {};
    \node[labeltext] (out#1) at (12.55,#2) {$y$};
    \node[wideblock] (m#1) at (7.65,#2-1.02) {#4};

    \draw[signal] (r#1) -- (s#1);
    \draw[signal] (s#1) -- (e#1);
    \draw[signal] (e#1) -- (c#1);
    \draw[signal] (c#1) -- (a#1);
    \draw[signal] (a#1) -- (p#1);
    \draw[signal] (p#1) -- (tap#1) -- (out#1);
    \draw[feedback] (tap#1) |- (m#1.east);
    \draw[feedback] (m#1.west) -| node[labeltext, left, pos=0.94] {$-$} (s#1.south);

    \node[plotbox] (pr#1) at (0,#2+0.68) {};
    \node[plotbox] (pe#1) at (3.20,#2+0.68) {};
    \node[plotbox] (pc#1) at (5.25,#2+0.68) {};
    \node[plotbox] (pa#1) at (7.65,#2+0.68) {};
    \node[plotbox] (py#1) at (10.15,#2+0.68) {};
    \node[plotbox, minimum width=2.05cm] (pm#1) at (7.65,#2-1.72) {};
  }

  \drawrow{A}{1.15}{无传感器滤波：噪声直接进入误差与控制量}{传感器}
  \drawrow{B}{-2.05}{加入一阶传感器滤波：反馈更平滑，执行器动作减轻}{传感器/滤波}
\end{tikzpicture}
\end{document}
